% GENERATED FILE. Edit canonical lessons, then run npm run generate:books.
% canonical-source-hash: fnv1a64:1346b912533fb64d
% canonical-lessons: PT-C281-facto, PT-C281-relatar, PT-C281-relato, PT-C281-relatorio, PT-C281-noticiar

\chapter{Os factos}
\label{ch:pt-os-factos}
\hlchaptermodality{281}

\begin{chapteropening}
\textbf{By the end of this chapter:} \emph{I can report the facts, give an account of what happened, and say that the news reported it.}

Complete the last lesson of chapter 281: I can report the facts, give an account of what happened, and say that the news reported it.
\end{chapteropening}

\section[o facto]{o facto — the fact}

\label{lesson:PT-C281-facto}



\begin{quote}
Before the new one: say the Portuguese for the edition, then the Portuguese for the article.
\end{quote}

\subsection*{You'll want to know: o facto}
\textbf{o facto} — \textquotedblleft{}the fact\textquotedblright{}.

\begin{quote}
É um facto: ninguém viu nada.
\end{quote}

\begin{quote}
\emph{It is a fact: nobody saw anything.}
\end{quote}

In Brazil: \emph{o fato}, which in Portugal is a suit. \emph{De facto} — in fact.

\begin{cousinweb}
\emph{facto} is a learned borrowing from Latin \emph{factum} \textquotedblleft{}a deed, a thing done\textquotedblright{} (EtymDB), the past participle of \emph{facere} \textquotedblleft{}to do, to make\textquotedblright{}; the inherited form of the same word is \emph{feito}.
\end{cousinweb}

\begin{grammarlens}[title={What you've built}]
One more word for telling, explaining and describing at B1.
\end{grammarlens}

\subsection*{Your turn}
\begin{itemize}
\raggedright
  \item \emph{Say it:} \emph{o facto}
  \item \emph{Say it:} \emph{o facto}, once more
  \item \emph{Say it:} a sentence of your own with \emph{o facto}
  \item \emph{Recall:} say the Portuguese for the edition, then the Portuguese for the article, then say \emph{o facto} again
\end{itemize}

\begin{tcolorbox}[breakable,colback=teal!4,colframe=teal!35!black,title={Before you move on}]
What does o facto mean? (\textquotedblleft{}The fact\textquotedblright{}.) Say it once more.
\end{tcolorbox}
\section[relatar]{relatar — to report, to give an account of}

\label{lesson:PT-C281-relatar}



\begin{quote}
Before the new one: say the Portuguese for the article, then the Portuguese for the fact.
\end{quote}

\subsection*{You'll want to know: relatar}
\textbf{relatar} — \textquotedblleft{}to report, to give an account of\textquotedblright{}.

\begin{quote}
O jornalista relatou tudo o que viu.
\end{quote}

\begin{quote}
\emph{The journalist reported everything he saw.}
\end{quote}

To tell in order what happened.

\begin{cousinweb}
\emph{relatar} is built on \emph{relato}, from Latin \emph{relātus} \textquotedblleft{}a telling\textquotedblright{} (EtymDB), the past participle of \emph{referre} \textquotedblleft{}to bring back, to report\textquotedblright{} (Lewis \& Short).
\end{cousinweb}

\begin{grammarlens}[title={What you've built}]
One more word for telling, explaining and describing at B1.
\end{grammarlens}

\subsection*{Your turn}
\begin{itemize}
\raggedright
  \item \emph{Say it:} \emph{relatar}
  \item \emph{Say it:} \emph{relatar}, once more
  \item \emph{Say it:} a sentence of your own with \emph{relatar}
  \item \emph{Recall:} say the Portuguese for the article, then the Portuguese for the fact, then say \emph{relatar} again
\end{itemize}

\begin{tcolorbox}[breakable,colback=teal!4,colframe=teal!35!black,title={Before you move on}]
What does relatar mean? (\textquotedblleft{}To report, to give an account of\textquotedblright{}.) Say it once more.
\end{tcolorbox}
\section[o relato]{o relato — the account, the narration}

\label{lesson:PT-C281-relato}



\begin{quote}
Before the new one: say the Portuguese for the fact, then the Portuguese for to report.
\end{quote}

\subsection*{You'll want to know: o relato}
\textbf{o relato} — \textquotedblleft{}the account, the narration\textquotedblright{}.

\begin{quote}
O relato dela foi muito claro.
\end{quote}

\begin{quote}
\emph{Her account was very clear.}
\end{quote}

A spoken or written telling of events.

\begin{cousinweb}
\emph{relato} is a learned borrowing from Latin \emph{relātus} \textquotedblleft{}a telling, a narration\textquotedblright{} (EtymDB), from \emph{referre} \textquotedblleft{}to bring back, to report\textquotedblright{}.
\end{cousinweb}

\begin{grammarlens}[title={What you've built}]
One more word for telling, explaining and describing at B1.
\end{grammarlens}

\subsection*{Your turn}
\begin{itemize}
\raggedright
  \item \emph{Say it:} \emph{o relato}
  \item \emph{Say it:} \emph{o relato}, once more
  \item \emph{Say it:} a sentence of your own with \emph{o relato}
  \item \emph{Recall:} say the Portuguese for the fact, then the Portuguese for to report, then say \emph{o relato} again
\end{itemize}

\begin{tcolorbox}[breakable,colback=teal!4,colframe=teal!35!black,title={Before you move on}]
What does o relato mean? (\textquotedblleft{}The account, the narration\textquotedblright{}.) Say it once more.
\end{tcolorbox}
\section[o relatório]{o relatório — the written report}

\label{lesson:PT-C281-relatorio}



\begin{quote}
Before the new one: say the Portuguese for to report, then the Portuguese for the account.
\end{quote}

\subsection*{You'll want to know: o relatório}
\textbf{o relatório} — \textquotedblleft{}the written report\textquotedblright{}.

\begin{quote}
O chefe pediu um relatório até sexta-feira.
\end{quote}

\begin{quote}
\emph{The boss asked for a report by Friday.}
\end{quote}

A formal document at work; \emph{a reportagem} belongs to the news.

\begin{cousinweb}
\emph{relatório} is built on \emph{relato} \textquotedblleft{}an account\textquotedblright{} with the ending \emph{-ório} (EtymDB), and goes back to Latin \emph{relātus} \textquotedblleft{}a telling\textquotedblright{}.
\end{cousinweb}

\begin{grammarlens}[title={What you've built}]
One more word for telling, explaining and describing at B1.
\end{grammarlens}

\subsection*{Your turn}
\begin{itemize}
\raggedright
  \item \emph{Say it:} \emph{o relatório}
  \item \emph{Say it:} \emph{o relatório}, once more
  \item \emph{Say it:} a sentence of your own with \emph{o relatório}
  \item \emph{Recall:} say the Portuguese for to report, then the Portuguese for the account, then say \emph{o relatório} again
\end{itemize}

\begin{tcolorbox}[breakable,colback=teal!4,colframe=teal!35!black,title={Before you move on}]
What does o relatório mean? (\textquotedblleft{}The written report\textquotedblright{}.) Say it once more.
\end{tcolorbox}
\section[noticiar]{noticiar — to announce in the news, to report}

\label{lesson:PT-C281-noticiar}



\begin{quote}
Before the new one: say the Portuguese for the account, then the Portuguese for the written report.
\end{quote}

\subsection*{You'll want to know: noticiar}
\textbf{noticiar} — \textquotedblleft{}to announce in the news, to report\textquotedblright{}.

\begin{quote}
A televisão noticiou a chegada do presidente.
\end{quote}

\begin{quote}
\emph{Television reported the president’s arrival.}
\end{quote}

What the media do with an event.

\begin{cousinweb}
\emph{noticiar} is built on \emph{notícia} \textquotedblleft{}news\textquotedblright{}, from Latin \emph{nōtitia} \textquotedblleft{}knowledge, a notice\textquotedblright{} (Lewis \& Short).
\end{cousinweb}

\begin{grammarlens}[title={What you've built}]
One more word for telling, explaining and describing at B1.
\end{grammarlens}

\subsection*{Your turn}
\begin{itemize}
\raggedright
  \item \emph{Say it:} \emph{noticiar}
  \item \emph{Say it:} \emph{noticiar}, once more
  \item \emph{Say it:} a sentence of your own with \emph{noticiar}
  \item \emph{Recall:} say the Portuguese for the account, then the Portuguese for the written report, then say \emph{noticiar} again
\end{itemize}

\begin{tcolorbox}[breakable,colback=teal!4,colframe=teal!35!black,title={Before you move on}]
What does noticiar mean? (\textquotedblleft{}To announce in the news, to report\textquotedblright{}.) Say it once more.
\end{tcolorbox}
